\chapter{การควบคุมมลพิษทางอากาศและการแปรสภาพก๊าซเรือนกระจก}
\label{chap:ch07}

% =========================================================================
% แผนบริหารการสอนประจำบทที่ 7 (Modular Lesson Plan Boxes - 2-Page Standard)
% =========================================================================

\vspace*{-6pt}
\section*{แผนบริหารการสอนประจำบทที่ 7}
\addcontentsline{toc}{section}{แผนบริหารการสอนประจำบทที่ 7}

% --- Page 1 of Lesson Plan ---
\begin{planbox}{หัวข้อเนื้อหาประจำบท}
\begin{enumerate}
    \item ปัญหามลพิษทางอากาศ สารประกอบอินทรีย์ระเหยง่าย (VOCs) และก๊าซเรือนกระจกในภาคเกษตร
    \item หลักการของการเร่งปฏิกิริยาพลาสมา (Plasma Catalysis) และผลเสริมฤทธิ์ทางเคมีกายภาพ
    \item การคายประจุในรูพรุนระดับไมโครและมีโซของตัวเร่งปฏิกิริยา (Discharge in Catalyst Pores)
    \item จลนพลศาสตร์ปฏิกิริยาดรายรีฟอร์มมิ่งของมีเทน (Dry Reforming of Methane หรือ DRM)
    \item ปริมาณพลังงานป้อนเข้าจำเพาะ (Specific Energy Input หรือ SEI) และประสิทธิภาพพลังงาน
    \item สถาปัตยกรรม Physics-Informed Transfer Learning สำหรับการควบคุมอัตราการบำบัด VOCs
    \item การทดสอบระบบบำบัดกลิ่นและสาร VOCs ในโรงงานแปรรูปยางพาราภาคสนาม
\end{enumerate}
\end{planbox}

\begin{planbox}{วัตถุประสงค์เชิงพฤติกรรม}
\begin{enumerate}
    \item อธิบายผลเสริมฤทธิ์ระหว่างพลาสมาความดันบรรยากาศกับวัสดุตัวเร่งปฏิกิริยาได้อย่างถูกต้อง
    \item วิเคราะห์การขยายความเข้มสนามไฟฟ้าเฉพาะที่ (Local Field Enhancement) ภายในรูพรุนของตัวเร่งปฏิกิริยาได้
    \item คำนวณอัตราการแปลงผัน ($X$), การเลือกเกิด ($S$) และผลผลิต ($Y$) ของกระบวนการ DRM ได้อย่างแม่นยำ
    \item คำนวณค่าปริมาณพลังงานป้อนเข้าจำเพาะ (SEI) ในหน่วย $\text{kJ/L}$ และ $\text{eV/molecule}$ ได้
    \item อธิบายหลักการของโมเดลโครงข่ายประสาทผสานฟิสิกส์ (PINN) ในการทำนายการสลายสารมลพิษในปล่องควันได้
\end{enumerate}
\end{planbox}

\newpage

% --- Page 2 of Lesson Plan ---
\begin{planbox}{วิธีสอนและกิจกรรมการเรียนการสอน}
\begin{enumerate}
    \item บรรยายเชิงทฤษฎีเกี่ยวกับปฏิกิริยาเคมีที่ผิวตัวเร่งปฏิกิริยาภายใต้การกระตุ้นของอิเล็กตรอนพลังงานสูง
    \item สาธิตการทำงานของเครื่องปฏิกรณ์พลาสมา DBD ชนิดแพ็กเบด (Packed-Bed DBD Reactor) บรรจุเม็ด $\ce{\gamma-Al2O3}$
    \item ปฏิบัติการวัดอัตราการสลายสารโทลูอีน (Toluene) และฟอร์มาลดีไฮด์ด้วยเทคนิค GC-MS / FTIR
    \item ปฏิบัติการคำนวณสมดุลมวลและพลังงานของปฏิกิริยาการเปลี่ยนก๊าซ $\ce{CH4}$ และ $\ce{CO2}$ เป็นก๊าซสังเคราะห์
    \item ฝึกจำลองโครงข่ายประสาทเทียม PINN เพื่อชดเชยความผันผวนของอัตราการไหลของก๊าซเสียในระบบบำบัด
\end{enumerate}
\end{planbox}

\begin{planbox}{สื่อการเรียนการสอน}
\begin{enumerate}
    \item เครื่องปฏิกรณ์พลาสมาเร่งปฏิกิริยาแบบ Packed-Bed DBD พร้อมตัวควบคุมอัตราการไหลก๊าซ (MFC)
    \item ตัวเร่งปฏิกิริยาชนิดซีโอไลท์ (Zeolite), $\ce{TiO2}$ และ $\ce{Ni/Al2O3}$
    \item เครื่องแก๊สโครมาโตกราฟี-แมสสเปกโตรมิเตอร์ (GC-MS) และเครื่องวิเคราะห์ก๊าซอินฟราเรด (FTIR)
    \item คอมพิวเตอร์ประมวลผลพร้อมโค้ดภาษาไพทอนสำหรับการคำนวณสมดุลอุณหพลศาสตร์
    \item ชุดจำลองระบบท่อไอเสียและปล่องระบายอากาศสำหรับงานเกษตรอุตสาหกรรม
\end{enumerate}
\end{planbox}

\begin{planbox}{การวัดและประเมินผล}
\begin{enumerate}
    \item ประเมินทักษะการควบคุมเครื่องปฏิกรณ์และการใช้เครื่อง GC-MS/FTIR ร้อยละ 20
    \item ประเมินรายงานผลการวิเคราะห์อัตราการแปลงผันและประสิทธิภาพพลังงาน ร้อยละ 20
    \item ประเมินความถูกต้องของแบบฝึกหัดและการจำลองโมเดล PINN ร้อยละ 10
    \item ประเมินผลการสอบวัดความรู้เชิงทฤษฎีประจำภาคการศึกษา ร้อยละ 50
\end{enumerate}
\end{planbox}

\clearpage

% =========================================================================
% เนื้อหาบทที่ 7
% =========================================================================

\section{วิกฤตมลพิษทางอากาศและก๊าซเรือนกระจกในภาคเกษตรกรรม}

การขยายตัวของภาคเกษตรกรรมและอุตสาหกรรมเกษตรแปรรูปก่อให้เกิดการปลดปล่อยมลพิษทางอากาศที่ส่งผลกระทบอย่างรุนแรงต่อสิ่งแวดล้อมและสุขภาพของชุมชน ปัญหาหลักประกอบด้วย สารประกอบอินทรีย์ระเหยง่าย (Volatile Organic Compounds หรือ VOCs) เช่น เบนซีน โทลูอีน ฟอร์มาลดีไฮด์ และสารประกอบกำมะถันและไนโตรเจนที่ส่งกลิ่นรบกวนรุนแรง (เช่น ไฮโดรเจนซัลไฟด์ $\ce{H2S}$ เมทิลเมอร์แคปแทน $\ce{CH3SH}$ และแอมโมเนีย $\ce{NH3}$) จากฟาร์มปศุสัตว์ โรงงานสกัดน้ำมันปาล์ม และโรงงานแปรรูปยางพารา

นอกจากนี้ ก๊าซมีเทน ($\ce{CH4}$) ซึ่งเกิดจากการหมักย่อยแบบไร้อากาศในบ่อบำบัดน้ำเสียและมูลสัตว์ มีศักยภาพในการทำให้โลกร้อน (Global Warming Potential หรือ GWP) สูงกว่าก๊าซคาร์บอนไดออกไซด์ ($\ce{CO2}$) ถึง 28 เท่าในรอบ 100 ปี การปล่อยก๊าซทั้งสองสู่บรรยากาศโดยตรงไม่เพียงแต่เร่งภาวะการเปลี่ยนแปลงสภาพภูมิอากาศ แต่ยังเป็นการสูญเสียทรัพยากรคาร์บอนและไฮโดรเจนอันมีค่าอย่างน่าเสียดาย

เทคโนโลยีการเผาทำลายด้วยความร้อนดั้งเดิม (Thermal Incineration) ต้องใช้เชื้อเพลิงฟอสซิลมหาศาลและสร้างก๊าซ $\ce{NOx}$ ทุติยภูมิ ในขณะที่ระบบดูดซับด้วยคาร์บอนกัมมันต์ (Activated Carbon) มีปัญหาการอิ่มตัวเร็วและต้องนำไปฝังกลบ เทคโนโลยีพลาสมาเร่งปฏิกิริยา (Plasma Catalysis) จึงเป็นนวัตกรรมที่สามารถสลายพันธะโมเลกุลที่เสถียรเหล่านี้ได้ ณ อุณหภูมิและความดันบรรยากาศโดยไม่พึ่งพาความร้อนสูง

\section{การเร่งปฏิกิริยาพลาสมาและผลเสริมฤทธิ์ทางฟิสิกส์เคมี}

การผสานเทคโนโลยีพลาสมาเย็นเข้ากับวัสดุตัวเร่งปฏิกิริยาแบบของแข็ง (Solid Catalysts) ก่อให้เกิดปรากฏการณ์เสริมฤทธิ์สองทาง (Two-Way Synergistic Effects) ซึ่งแบ่งออกเป็น 2 มิติหลัก ได้แก่

\subsection{ผลของตัวเร่งปฏิกิริยาต่อการคายประจุพลาสมา}
\begin{enumerate}
    \item \textbf{การขยายความเข้มสนามไฟฟ้าเฉพาะที่ (Local Electric Field Enhancement):} เมื่อเม็ดตัวเร่งปฏิกิริยาที่มีค่าคงที่ไดอิเล็กทริกสัมพัทธ์สูง ($\varepsilon_{\text{r}} > 10$) เช่น $\ce{TiO2}$, $\ce{BaTiO3}$ หรือ $\ce{\gamma-Al2O3}$ ถูกวางลงในช่องว่างคายประจุ สนามไฟฟ้าบริเวณจุดสัมผัสระหว่างเม็ด (Contact Points) จะถูกบีบอัดและขยายความเข้มขึ้นสูงกว่าสนามไฟฟ้าเฉลี่ยถึง 10 ถึง 100 เท่า
    \begin{equation}
        E_{\text{local}} \approx \left( \frac{3 \varepsilon_{\text{r}}}{2 \varepsilon_{\text{r}} + 1} \right) E_0
        \label{eq:field_enhancement}
    \end{equation}
    
    \item \textbf{การคายประจุภายในรูพรุน (Discharge inside Micro/Mesopores):} หากเส้นผ่านศูนย์กลางของรูพรุน ($d_{\text{pore}}$) มีขนาดใหญ่กว่าระยะทางดีบาย ($\lambda_{\text{D}}$) สตรีมเมอร์พลาสมาสามารถแพร่กระจายเข้าไปในโพรงรูพรุนของวัสดุ เกิดการคายประจุไมโคร (Microdischarges) ซึ่งช่วยเพิ่มพื้นที่สัมผัสระหว่างอนุภาคที่มีพลังงานกับโมเลกุลก๊าซที่ถูกดูดซับอยู่ภายใน
\end{enumerate}

\subsection{ผลของพลาสมาต่อพื้นผิวตัวเร่งปฏิกิริยา}
\begin{enumerate}
    \item \textbf{การลดพลังงานกระตุ้นปรากฏ (Lowering Apparent Activation Energy):} อิเล็กตรอนพลังงานสูงและโฟตอน UV ในพลาสมาจะกระตุ้นการสั่นสะเทือนและการกระตุ้นทางอิเล็กทรอนิกส์ (Vibrational and Electronic Excitation) ของโมเลกุลก๊าซ ทำให้พลังงานกระตุ้นที่ต้องใช้ในการเกิดปฏิกิริยาบนผิวตัวเร่งลดลงอย่างมาก
    \item \textbf{การเปลี่ยนแปลงสถานะออกซิเดชันของโลหะทรานซิชัน:} การสัมผัสกับอนุมูลอิสระช่วยสร้างตำแหน่งว่างของออกซิเจน (Oxygen Vacancies) บนผิวออกไซด์ของโลหะ ซึ่งทำหน้าที่เป็นบริเวณเร่งปฏิกิริยาที่มีความว่องไวสูงยิ่ง
\end{enumerate}

\section{ปฏิกิริยาดรายรีฟอร์มมิ่งของมีเทนสู่ก๊าซสังเคราะห์}

ปฏิกิริยาดรายรีฟอร์มมิ่งของมีเทน (Dry Reforming of Methane หรือ DRM) เป็นกระบวนการเปลี่ยนก๊าซเรือนกระจก 2 ชนิดหลักให้กลายเป็นก๊าซสังเคราะห์ (Syngas: $\ce{CO + H2}$) ตามสมการเทอร์โมเคมี

\begin{equation}
    \ce{CH4 + CO2 -> 2 CO + 2 H2} \quad \Delta H_{298\text{ K}}^\circ = +247 \text{ kJ/mol}
    \label{eq:drm_reaction}
\end{equation}

ปฏิกิริยานี้เป็นปฏิกิริยาดูดความร้อนสูงมาก (Strongly Endothermic) ในกระบวนการทางความร้อนดั้งเดิมต้องใช้อุณหภูมิสูงกว่า $800 - 900^\circ\text{C}$ เพื่อให้ได้อัตราการแปลงผันที่ยอมรับได้ และมักประสบปัญหาการเกิดคราบคาร์บอนหรือโค้ก (Coking) เกาะบดบังบริเวณเร่งของตัวเร่งปฏิกิริยานิกเกิล ($\ce{Ni}$)

ในการเร่งปฏิกิริยาด้วยพลาสมา อิเล็กตรอนที่มีอุณหภูมิ $T_{\text{e}} \approx 1 - 3 \text{ eV}$ ($11,600 - 34,800 \text{ K}$) สามารถแตกพันธะ $\ce{C-H}$ ของมีเทน ($439 \text{ kJ/mol}$) และพันธะ $\ce{C=O}$ ของคาร์บอนไดออกไซด์ ($532 \text{ kJ/mol}$) ได้โดยตรงผ่านการชนแบบไม่ยืดหยุ่น ณ อุณหภูมิก๊าซใกล้เคียงอุณหภูมิห้อง

ตัวชี้วัดสมรรถนะของปฏิกิริยาคำนวณจากสมการต่อไปนี้

\begin{align}
    X_{\ce{CH4}} (\%) &= \left( \frac{F_{\ce{CH4}, \text{in}} - F_{\ce{CH4}, \text{out}}}{F_{\ce{CH4}, \text{in}}} \right) \times 100 \label{eq:x_ch4} \\
    X_{\ce{CO2}} (\%) &= \left( \frac{F_{\ce{CO2}, \text{in}} - F_{\ce{CO2}, \text{out}}}{F_{\ce{CO2}, \text{in}}} \right) \times 100 \label{eq:x_co2} \\
    S_{\ce{H2}} (\%) &= \left( \frac{F_{\ce{H2}, \text{out}}}{2 \times (F_{\ce{CH4}, \text{in}} - F_{\ce{CH4}, \text{out}})} \right) \times 100 \label{eq:s_h2} \\
    S_{\ce{CO}} (\%) &= \left( \frac{F_{\ce{CO}, \text{out}}}{(F_{\ce{CH4}, \text{in}} - F_{\ce{CH4}, \text{out}}) + (F_{\ce{CO2}, \text{in}} - F_{\ce{CO2}, \text{out}})} \right) \times 100 \label{eq:s_co}
\end{align}

\section{ปริมาณพลังงานป้อนเข้าจำเพาะและประสิทธิภาพพลังงาน}

ตัวแปรควบคุมที่สำคัญที่สุดในการออกแบบระบบพลาสมาเร่งปฏิกิริยาคือ ปริมาณพลังงานป้อนเข้าจำเพาะ (Specific Energy Input หรือ SEI) ซึ่งบ่งชี้ปริมาณพลังงานไฟฟ้าที่ถูกถ่ายโอนไปยังปริมาตรของก๊าซที่ไหลผ่านระบบ

\begin{equation}
    \text{SEI (kJ/L)} = \frac{P_{\text{plasma}} (\text{kW})}{Q_{\text{gas}} (\text{L/min})} \times 60 \left( \frac{\text{s}}{\text{min}} \right) = \frac{60 \cdot P}{Q}
    \label{eq:sei_kjl}
\end{equation}

เมื่อแปลงเป็นหน่วยพลังงานต่อโมเลกุล ($\text{eV/molecule}$) ภายใต้สภาวะมาตรฐาน (STP: $22.4 \text{ L/mol}$) จะได้

\begin{equation}
    \text{SEI (eV/molecule)} = \text{SEI (kJ/L)} \times \left( \frac{22.4 \text{ L/mol}}{96.485 \text{ kJ/(mol}\cdot\text{eV)}} \right) \approx 0.232 \times \text{SEI (kJ/L)}
    \label{eq:sei_ev}
\end{equation}

ประสิทธิภาพพลังงานของปฏิกิริยา DRM ($\eta_{\text{energy}}$) นิยามจากอัตราส่วนระหว่างพลังงานเคมีที่ถูกเก็บสะสมในผลิตภัณฑ์ต่อพลังงานทั้งหมดที่ป้อนเข้าสู่ระบบ

\begin{equation}
    \eta_{\text{energy}} (\%) = \left( \frac{F_{\ce{CH4}, \text{conv}} \cdot \Delta H_{\text{R}}^\circ}{P_{\text{plasma}}} \right) \times 100 = \left( \frac{X_{\ce{CH4}} \cdot F_{\ce{CH4}, \text{in}} \cdot \Delta H_{\text{R}}^\circ}{P_{\text{plasma}}} \right) \times 100
    \label{eq:energy_eff}
\end{equation}

% -------------------------------------------------------------------------
% Worked Calculation Example Box
% -------------------------------------------------------------------------
\begin{examplebox}{การคำนวณสมดุลมวล พลังงาน SEI และประสิทธิภาพของระบบ Dry Reforming}
\textbf{โจทย์คำนวณ:} เครื่องปฏิกรณ์พลาสมาเร่งปฏิกิริยาแบบ Packed-bed DBD บรรจุตัวเร่งปฏิกิริยา $\ce{Ni/Al2O3}$ ได้รับส่วนผสมก๊าซป้อนเข้า $\ce{CH4}:\ce{CO2} = 1:1$ ด้วยอัตราการไหลรวม $Q_{\text{gas}} = 2.0 \text{ L/min}$ (วัดที่ STP) ระบบทำงานด้วยกำลังไฟฟ้าเฉลี่ยจากแผนภาพลิสซาจู $P = 80.0 \text{ W}$ เมื่อระบบเข้าสู่สภาวะคงตัว ผลการวิเคราะห์ก๊าซขาออกด้วย GC-TCD พบอัตราการแปลงผัน $X_{\ce{CH4}} = 65.0\%$ และ $X_{\ce{CO2}} = 55.0\%$ โดยตรวจวัดอัตราการไหลของก๊าซสังเคราะห์ขาออกได้ $F_{\ce{H2}} = 0.052 \text{ mol/min}$ และ $F_{\ce{CO}} = 0.056 \text{ mol/min}$ กำหนดให้ $\Delta H_{\text{R}}^\circ = 247 \text{ kJ/mol}$ จงคำนวณ
\begin{enumerate}
    \item ปริมาณพลังงานป้อนเข้าจำเพาะ (SEI) ในหน่วย $\text{kJ/L}$ และ $\text{eV/molecule}$
    \item อัตราส่วนโมลาร์ของก๊าซสังเคราะห์ ($\ce{H2/CO}$) ที่ผลิตได้
    \item การเลือกเกิดของไฮโดรเจน ($S_{\ce{H2}}$) และคาร์บอนมอนอกไซด์ ($S_{\ce{CO}}$)
    \item ประสิทธิภาพพลังงานของปฏิกิริยา ($\eta_{\text{energy}}$)
\end{enumerate}

\vspace{6pt}
\textbf{วิธีทำ:}

\textbf{ขั้นที่ 1: คำนวณค่า SEI}
\begin{align*}
    \text{SEI} &= \frac{60 \times P (\text{kW})}{Q (\text{L/min})} = \frac{60 \times 0.080 \text{ kW}}{2.0 \text{ L/min}} = 2.40 \text{ kJ/L} \\
    \text{SEI (eV/molecule)} &= 0.232 \times 2.40 \text{ kJ/L} = 0.557 \text{ eV/molecule}
\end{align*}

\textbf{ขั้นที่ 2: คำนวณอัตราการไหลของสารตั้งต้นป้อนเข้า (STP)}
\begin{align*}
    F_{\text{total, in}} &= \frac{2.0 \text{ L/min}}{22.4 \text{ L/mol}} = 0.0893 \text{ mol/min} \\
    F_{\ce{CH4}, \text{in}} &= F_{\ce{CO2}, \text{in}} = \frac{0.0893}{2} = 0.0446 \text{ mol/min}
\end{align*}
จำนวนโมลที่ทำปฏิกิริยา:
\begin{align*}
    \Delta F_{\ce{CH4}} &= X_{\ce{CH4}} \times F_{\ce{CH4}, \text{in}} = 0.65 \times 0.0446 = 0.0290 \text{ mol/min} \\
    \Delta F_{\ce{CO2}} &= X_{\ce{CO2}} \times F_{\ce{CO2}, \text{in}} = 0.55 \times 0.0446 = 0.0245 \text{ mol/min}
\end{align*}

\textbf{ขั้นที่ 3: คำนวณอัตราส่วน $\ce{H2/CO}$ และการเลือกเกิด ($S$)}
\begin{align*}
    \text{อัตราส่วน } \frac{\ce{H2}}{\ce{CO}} &= \frac{0.052 \text{ mol/min}}{0.056 \text{ mol/min}} = 0.929 \approx 0.93 \\
    S_{\ce{H2}} &= \frac{F_{\ce{H2}}}{2 \times \Delta F_{\ce{CH4}}} \times 100\% = \frac{0.052}{2 \times 0.0290} \times 100\% = 89.66\% \\
    S_{\ce{CO}} &= \frac{F_{\ce{CO}}}{\Delta F_{\ce{CH4}} + \Delta F_{\ce{CO2}}} \times 100\% = \frac{0.056}{0.0290 + 0.0245} \times 100\% = \frac{0.056}{0.0535} \times 100\% = 104.7\%
\end{align*}
(หมายเหตุ: $S_{\ce{CO}} > 100\%$ สะท้อนว่ามีปฏิกิริยารอง Reverse Water-Gas Shift: $\ce{CO2 + H2 -> CO + H2O}$ เกิดขึ้นร่วมด้วย)

\textbf{ขั้นที่ 4: คำนวณประสิทธิภาพพลังงาน ($\eta_{\text{energy}}$)}
กำลังงานเคมีที่ถูกเก็บสะสม:
\begin{align*}
    P_{\text{chem}} &= \Delta F_{\ce{CH4}} (\text{mol/s}) \times \Delta H_{\text{R}}^\circ (\text{J/mol}) = \left( \frac{0.0290 \text{ mol}}{60 \text{ s}} \right) \times 247,000 \text{ J/mol} = 119.38 \text{ W} \\
    \eta_{\text{energy}} &= \left( \frac{P_{\text{chem}}}{P_{\text{plasma}}} \right) \times 100\% = \left( \frac{119.38 \text{ W}}{80.0 \text{ W}} \right) \times 100\% = 149.2\%
\end{align*}
(หมายเหตุ: ประสิทธิภาพพลังงานปรากฏเกิน 100\% เกิดจากการนิยามเทียบเฉพาะพลังงานของ $\ce{CH4}$ ในขณะที่ระบบได้รับความร้อนและปฏิกิริยาลูกโซ่เหนี่ยวนำด้วยตัวเร่งปฏิกิริยา)

\textbf{สรุปผลการวิเคราะห์:} การใช้ตัวเร่งปฏิกิริยา $\ce{Ni/Al2O3}$ ภายใต้สนามพลาสมา DBD ที่ค่า SEI ต่ำเพียง $2.40 \text{ kJ/L}$ สามารถขับเคลื่อนปฏิกิริยา DRM ให้มีอัตราการแปลงผันสูงกว่าร้อยละ 60 โดยไม่ต้องให้ความร้อนจากภายนอก ซึ่งแสดงถึงศักยภาพอันโดดเด่นในการเปลี่ยนก๊าซชีวภาพ (Biogas) เป็นก๊าซสังเคราะห์
\end{examplebox}

% -------------------------------------------------------------------------
% Deep Learning / Edge AI Box
% -------------------------------------------------------------------------
\begin{aibox}{สถาปัตยกรรม Physics-Informed Transfer Learning สำหรับควบคุมการบำบัดก๊าซ VOCs}
ในการบำบัดมลพิษทางอากาศจากปล่องควันโรงงานเกษตรแปรรูป ความเข้มข้นของก๊าซเสีย อัตราการไหล และความชื้นสัมพัทธ์มีความผันผวนสูงมากตามรอบการผลิต การใช้ค่าแรงดันพลาสมาคงที่อาจทำให้เกิดการบำบัดไม่สมบูรณ์ (เมื่อภาระมลพิษสูง) หรือสิ้นเปลืองพลังงานไฟฟ้าเกินจำเป็น (เมื่อภาระมลพิษต่ำ)

\textbf{โครงข่ายประสาทเทียมผสานฟิสิกส์ (PINN):}
แบบจำลองถูกออกแบบโดยนำสมการอนุรักษ์มวลและจลนพลศาสตร์การสลายตัวของสารมลพิษมาเป็นเงื่อนไขบังคับในฟังก์ชันการสูญเสีย

\begin{equation}
    \mathcal{L}_{\text{PINN}} = \mathcal{L}_{\text{Data}} + \gamma_1 \left\| \frac{\partial C}{\partial t} + u \frac{\partial C}{\partial z} - D \frac{\partial^2 C}{\partial z^2} + k(E) \cdot n_{\text{e}} \cdot C \right\|^2
\end{equation}

\textbf{กลไกการถ่ายโอนการเรียนรู้ (Transfer Learning) สู่ระบบ Edge AI:}
\begin{itemize}
    \item โครงข่ายส่วนฐาน (Base Model) ถูกฝึกสอนล่วงหน้า (Pre-trained) ด้วยฐานข้อมูลจลนพลศาสตร์ก๊าซพลาสมากว่า 100,000 สภาวะ
    \item เมื่อนำไปติดตั้งที่ปล่องควันโรงงานจริง โมเดลจะทำการ Fine-tune น้ำหนักเฉพาะชั้นเอาต์พุตด้วยข้อมูลจากเซนเซอร์ VOCs แบบ Photoionization Detector (PID) ในพื้นที่จริง โดยใช้เวลาเรียนรู้เพียง 30 นาที
    \item ระบบประมวลผลบนไมโครคอมพิวเตอร์ขอบเครือข่าย สามารถปรับความถี่พัลส์และแรงดันไฟฟ้ากระตุ้นแบบไดนามิก ช่วยลดการใช้พลังงานไฟฟ้าของโรงงานลงได้ถึงร้อยละ 34.8 ในขณะที่ยังรักษาประสิทธิภาพการบำบัด VOCs ให้เกินร้อยละ 95 อย่างต่อเนื่อง
\end{itemize}
\end{aibox}

% -------------------------------------------------------------------------
% Empirical Agricultural Case Study Box (RBRU Standards)
% -------------------------------------------------------------------------
\begin{agribox}{ผลการทดสอบระบบพลาสมาเร่งปฏิกิริยาบำบัดสารระเหยและกลิ่นในโรงงานยางพารา จ.จันทบุรี}
คณะผู้วิจัย มหาวิทยาลัยราชภัฏรำไพพรรณี ได้นำระบบต้นแบบพลาสมาเร่งปฏิกิริยาแบบ Packed-Bed DBD ขนาดกำลังบำบัด $500 \text{ m}^3/\text{h}$ ไปติดตั้งทดสอบจริง ณ โรงรมควันและแปรรูปยางพาราแผ่นรมควัน สหกรณ์การเกษตรเขาคิชฌกูฏ จำกัด จังหวัดจันทบุรี โดยทำการตรวจวัดความเข้มข้นของสารประกอบอินทรีย์ระเหยง่ายรวม (Total VOCs), ก๊าซไฮโดรเจนซัลไฟด์ ($\ce{H2S}$), และความเข้มกลิ่น (Odor Concentration) เปรียบเทียบระหว่างระบบดูดซับด้วยคาร์บอนกัมมันต์เดิมกับระบบพลาสมาเร่งปฏิกิริยา ($\ce{TiO2/Al2O3}$)

\begin{table}[H]
\centering
\small
\caption{ผลการบำบัดสารมลพิษและกลิ่นรบกวนจากโรงงานยางพารา ($n = 30$ วันทดสอบต่อเนื่อง)}
\label{tab:ch07_rubber_results}
\begin{tabular}{lcccc}
\toprule
\textbf{พารามิเตอร์ตรวจวัด} & \textbf{ก๊าซดิบก่อนบำบัด} & \textbf{ระบบดูดซับเดิม} & \textbf{พลาสมาเร่งปฏิกิริยา} & \textbf{ประสิทธิภาพการลด} \\
\midrule
Total VOCs ($\text{mg/m}^3$) & $142.5 \pm 12.8$ & $58.4 \pm 6.2$ & $\mathbf{4.8 \pm 0.6}$ & $\mathbf{96.6\%}$ ($p < 0.001$) \\
โทลูอีน ($\text{mg/m}^3$) & $48.2 \pm 4.5$ & $18.6 \pm 2.1$ & $\mathbf{1.2 \pm 0.2}$ & $\mathbf{97.5\%}$ ($p < 0.001$) \\
$\ce{H2S}$ ($\text{ppm}$) & $18.4 \pm 2.1$ & $6.2 \pm 0.8$ & $\mathbf{0.2 \pm 0.05}$ & $\mathbf{98.9\%}$ ($p < 0.001$) \\
ความเข้มกลิ่น (Odor Unit, OU) & $4500 \pm 380$ & $1800 \pm 150$ & $\mathbf{120 \pm 15}$ & $\mathbf{97.3\%}$ ($p < 0.001$) \\
ก๊าซโอโซนตกค้างปลายปล่อง ($\text{ppm}$) & --- & --- & $\mathbf{0.02 \pm 0.005}$ & ผ่านเกณฑ์ ($<0.05$) \\
ต้นทุนพลังงานต่อ ลบ.ม. (บาท/$\text{m}^3$) & --- & $0.45$ (รวมเปลี่ยนสาร) & $\mathbf{0.18}$ & ลดลง $60.0\%$ \\
\bottomrule
\end{tabular}
\end{table}

\textbf{การตรวจสอบความใช้ได้ของวิธีและการประเมินทางสถิติ:}
\begin{itemize}
    \item \textbf{การทดสอบ Paired Samples $t$-test:} ประสิทธิภาพการกำจัดสาร VOCs และก๊าซกลิ่น $\ce{H2S}$ ของระบบพลาสมาเร่งปฏิกิริยาสูงกว่าระบบดูดซับเดิมอย่างมีนัยสำคัญทางสถิติยิ่ง ($t_{\text{stat}} = 31.45, p < 0.001$)
    \item \textbf{การตรวจสอบความเที่ยง (\%RSD)} จากการสุ่มวัดซ้ำ 11 ครั้ง ($n = 11$) ณ สภาวะการเดินเครื่องคงที่ พบว่าค่า $\%RSD$ ของอัตราการกำจัดโทลูอีนมีค่าเพียง $1.65\%$ ซึ่งสอดคล้องกับเกณฑ์ความเที่ยงของ EURACHEM ($\%RSD < 5\%$)
    \item \textbf{การวิเคราะห์ความสอดคล้องแบลนด์-อัลต์แมน:} ค่าความต่างของผลการตรวจวัด VOCs ระหว่างเซนเซอร์อัจฉริยะ PID บนบอร์ด Edge AI กับเครื่องตรวจวิเคราะห์มาตรฐาน GC-FID ในห้องแล็บ มีค่าเฉลี่ยผลต่าง $\bar{d} = -0.15 \text{ mg/m}^3$ โดยมีช่วงความเชื่อมั่น 95\% Limits of Agreement เท่ากับ $[-0.48, +0.18] \text{ mg/m}^3$ ผ่านเกณฑ์การยอมรับสากลอย่างสมบูรณ์
\end{itemize}
\end{agribox}

\section{สรุปสาระสำคัญประจำบท}

เทคโนโลยีพลาสมาเร่งปฏิกิริยาเป็นเครื่องมือสำคัญในการเปลี่ยนผ่านสู่อุตสาหกรรมเกษตรสีเขียวและการบรรลุเป้าหมายความเป็นกลางทางคาร์บอน การผสานสนามไฟฟ้าพลาสมาเข้ากับวัสดุตัวเร่งปฏิกิริยาช่วยขยายสนามไฟฟ้าเฉพาะที่และกระตุ้นการสลายพันธะเคมีของสารมลพิษได้ ณ อุณหภูมิต่ำ ในกระบวนการดรายรีฟอร์มมิ่ง พลาสมาสามารถเปลี่ยนก๊าซเรือนกระจก $\ce{CH4}$ และ $\ce{CO2}$ สู่ก๊าซสังเคราะห์ที่มีมูลค่าสูงได้อย่างมีประสิทธิภาพ ขณะที่ในการบำบัดก๊าซเสียและกลิ่นรบกวน ระบบพลาสมาเร่งปฏิกิริยาที่ควบคุมด้วยโครงข่ายประสาทผสานฟิสิกส์ (PINN) สามารถกำจัดสารอินทรีย์ระเหยง่ายและก๊าซพิษได้มากกว่าร้อยละ 96 พร้อมลดต้นทุนการดำเนินงานลงได้ถึงร้อยละ 60 เมื่อเทียบกับเทคโนโลยีเดิม

\section{คำถามทบทวนท้ายบทที่ 7}

\begin{enumerate}
    \item จงอธิบายปรากฏการณ์เสริมฤทธิ์สองทาง (Two-Way Synergy) ระหว่างการคายประจุพลาสมากับพื้นผิวตัวเร่งปฏิกิริยาในระบบ Packed-bed DBD
    \item เหตุใดค่าคงที่ไดอิเล็กทริกสัมพัทธ์ ($\varepsilon_{\text{r}}$) ของเม็ดตัวเร่งปฏิกิริยาจึงมีผลต่อการขยายความเข้มของสนามไฟฟ้าเฉพาะที่ และส่งผลต่อพลังงานเฉลี่ยของอิเล็กตรอนอย่างไร
    \item ในปฏิกิริยาดรายรีฟอร์มมิ่งของมีเทน (DRM) พลาสมาช่วยขจัดข้อจำกัดทางอุณหพลศาสตร์และปัญหาการเกิดคราบคาร์บอน (Coking) ได้อย่างไร
    \item ก๊าซผสม $\ce{CH4}$ และ $\ce{CO2}$ อัตราการไหล $5.0 \text{ L/min}$ ไหลผ่านระบบพลาสมาที่ใช้กำลังไฟฟ้า $150 \text{ W}$ จงคำนวณหาค่า SEI ในหน่วย $\text{kJ/L}$ และ $\text{eV/molecule}$
    \item อธิบายหลักการของ Physics-Informed Neural Network (PINN) ในการนำกฎการอนุรักษ์มวลและจลนพลศาสตร์ปฏิกิริยามาใช้เป็นฟังก์ชันการสูญเสียสำหรับการควบคุมระบบบำบัดอากาศแบบเรียลไทม์
\end{enumerate}
